\chapter{Circuitos recortadores con diodos zener}

\begin{figure}[H]
  \centering
  \begin{tikzpicture}
    \begin{axis}[
      width=12cm,
      height=8cm,
      xlabel={$t$ [ms]},
      ylabel={$V$ [V]},
      grid=both,
      minor tick num=1,
      scaled ticks=false,
      xmax=40, xmin=0,
      extra y ticks={-5.1, 20},
      extra y tick style={
        grid style={darkgray, dashed, opacity=0.8},
        tick style={darkgray},
        tick label style={darkgray}
       },
    ]
    \addplot[
      color=blue,
      mark=none,
      dashed,
    ] table[
      col sep=tab,
      header=true,
      x expr=\thisrow{time}*1000,
      y = V(n003)
    ] {tables/RecortadorVin.txt};
    \addlegendentry{$v_{in}$}
    \addplot[
      color=red,
      mark=none,
      thick,
    ] table[
      col sep=tab,
      header=true,
      x expr=\thisrow{time}*1000,
      y = V(n004)
    ] {tables/RecortadorVoutA.txt};
    \addlegendentry{$v_{out}$}
    \end{axis}
  \end{tikzpicture}
  \caption{Señal de salida del circuito 1, obtenida por simulación.}
  \label{fig:RecortadorA}
\end{figure}

\begin{figure}[H]
  \centering
  \begin{tikzpicture}
    \begin{axis}[
      width=12cm,
      height=8cm,
      xlabel={$t$ [ms]},
      ylabel={$V$ [V]},
      grid=both,
      minor tick num=1,
      scaled ticks=false,
      xmax=40, xmin=0,
      extra y ticks={-12, 6.8},
      extra y tick style={
        grid style={darkgray, dashed, opacity=0.8},
        tick style={darkgray},
        tick label style={darkgray}
       },
    ]
    \addplot[
      color=blue,
      mark=none,
      dashed,
    ] table[
      col sep=tab,
      header=true,
      x expr=\thisrow{time}*1000,
      y = V(n003)
    ] {tables/RecortadorVin.txt};
    \addlegendentry{$v_{in}$}
    \addplot[
      color=red,
      mark=none,
      thick,
    ] table[
      col sep=tab,
      header=true,
      x expr=\thisrow{time}*1000,
      y = V(n002)
    ] {tables/RecortadorVoutB.txt};
    \addlegendentry{$v_{out}$}
    \end{axis}
  \end{tikzpicture}
  \caption{Señal de salida del circuito 2, obtenida por simulación.}
  \label{fig:RecortadorB}
\end{figure}
